\chapter{Of Emergency and the Suspension of Ordinary Law}
\label{art:emergency}

\articlenote{In which is set down what constitutes an emergency that threatens the whole silo, and how the Mayor, Judicial, and the Heads of Departments may act beyond the ordinary Pact to preserve the silo's survival.}

\section{The Declaration of Emergency}

An emergency exists when the silo faces a threat to its survival or the survival of a substantial portion of its citizens: catastrophic damage to the water or air systems, failure of the generator, breach of the outer walls, disease threatening the whole silo, or such other emergency as the Mayor and Judicial jointly determine. In an emergency, ordinary procedures may be set aside, but only by the Mayor and Judicial acting together, and never for longer than one season without renewal of the emergency by Assembly vote.

\begin{clauses}
\item The Mayor may declare an emergency and notify Judicial and the Heads of Departments. Judicial may challenge the declaration, and if challenged, Judicial shall rule whether the emergency is genuine, and the ruling is final.
\item An emergency suspends this Pact's ordinary procedures for a time set by the Mayor and Judicial jointly, not to exceed one season; if the emergency continues beyond one season, the Mayor must petition the Assembly for renewal of the emergency, and the Assembly must ratify the renewal by simple majority vote.
\item During an emergency, the Mayor may order citizens to work beyond their ordinary labor hours, may redirect resources from one Department to another as the emergency requires, and may restrict citizen movement between floors as necessary to preserve the silo's essential functions.
\item During an emergency, Judicial may hear and judge an offense with such swiftness as the emergency requires, and the hearing Article 6 grants for an offense not among the gravest may be shortened but not refused. No measure of correction beyond those Article 17 names may be given, and no offense not named among the gravest may be answered by cleaning, in emergency as in ordinary times, whatever order is given by any office.
\end{clauses}

\section{The Green List and Emergency Protocol}

In advance of any known emergency, the Mayor and the Heads of Departments maintain a protocol, called the green list, known to Judicial and to the Head of every Department, setting down in advance what measures must be taken in each kind of emergency --- the sealing of certain floors, the rationing of water or air, the order in which the generator's power is given and withheld, the shutting down of what the silo can spare. This list is updated yearly and kept in the Judicial archive.

\begin{clauses}
\item A citizen ordered to follow the green list protocol during an emergency shall obey that order as though it were a standing order of the Mayor, and may not refuse or delay the order, save Judicial rules during or after the emergency that the order was not a genuine emergency requirement.
\item The green list is not part of this Pact and is not open to citizen review or publication, save where a citizen has appealed an order given under the green list and Judicial is reviewing whether the order was lawful. The green list is maintained in secure archive under Article 21 and its contents and updates are the sole determination of the Mayor and Heads of Departments.
\item The green list is maintained by the Mayor in consultation with Judicial and the Heads of every Department, and is updated not less than once yearly, before the Assembly's ordinary session; the fact of its updating, but not its contents, is reported to the Assembly at that session, so that the Heads of Departments know the current protocol and every citizen knows that one is kept.
\end{clauses}

\section{The Duration and the End of Emergency}

An emergency, once declared, holds for the time set by the Mayor and Judicial. At the end of that time, the emergency is lifted, and the Pact's ordinary procedures resume in full. If the emergency continues, it must be renewed.

\begin{clauses}
\item At the end of the time set for the emergency, all extraordinary measures cease, and all citizens released to other duties resume those duties. Any citizen held or assigned to emergency labor shall be reassigned to ordinary labor or compensated for the extraordinary work by the Mayor or the Department concerned.
\item An emergency renewed beyond one season requires Assembly ratification; an emergency lasting longer than two seasons must be ratified anew by Assembly vote once per season thereafter.
\item Any order given during an emergency that violates this Pact in a way not authorized by the emergency declaration is answerable under Article 17, and a citizen obeying such an order is not punished; the responsibility rests with the officer who gave the unlawful order.
\end{clauses}

\section{The Holding of a Final Request in Abeyance}

In the gravest emergency, where the silo's survival itself is at stake, the Mayor and Judicial may jointly hold in abeyance the requirement under Article 18 that a citizen requesting to go outside be sent to the cleaning without delay. A citizen whose request is so held remains in the silo and shall be assigned to such labor as the emergency requires.

\begin{clauses}
\item A citizen's final request under Article 18 may be held in abeyance, not suspended, during an emergency, by the Mayor and Judicial jointly, for the duration of the emergency only; a citizen whose request is held in abeyance may renew it at any time after the emergency ends.
\item A citizen whose request is held in abeyance may petition Judicial for a ruling that the emergency no longer poses a genuine threat to the silo's survival and that the citizen's request should be honored; Judicial shall rule within one cycle of the petition.
\item No citizen's request shall be permanently denied or forgotten; if the emergency lasts longer than the citizen's own life, the citizen's death is recorded in the rolls without the cleaning, and the request is entered in the archive as unfulfilled but not denied.
\end{clauses}

\section{The Necessity and the Restraint}

Emergency powers exist to preserve the silo; they are not given to expand the Mayor's or Judicial's authority. A citizen may petition Judicial after an emergency to rule whether any order given during the emergency was necessary to the survival, and Judicial may order compensation or restitution if the order was deemed unnecessary.

\begin{clauses}
\item After an emergency ends, Judicial shall, within one season, review any emergency orders given and determine whether they were necessary to the stated emergency purpose. Orders found to be unnecessary or excessive may trigger compensation to citizens harmed by them.
\item A citizen who dies or is injured in service to an emergency order determined by Judicial to have been necessary is remembered in the archive, and the citizen's household may petition for compensation under Article 17's restitution provisions.
\end{clauses}

\section{The Announcement of an Emergency}

An emergency binds no citizen who has not been told of it, and the telling is done by one road for every floor.

\begin{clauses}
\item Upon the declaration under Section 1, the Mayor sends the fact of it, the kind of emergency, and the measures first ordered by the stations of official word under Article 8 to the Sheriff, the Judge, and the Head of every Department, and by the same stations to each Deputy's station under Article 7; the Deputies send Peacekeepers to every landing of their zones, who cry the emergency at the landing and post the Mayor's notice there, in the Mayor's own words, within one hour of its reaching the station.
\item A citizen who hears the emergency cried goes to the citizen's own dwelling, or to the citizen's own floor of labor where the citizen's trade is one the emergency requires, and remains there until a further order reaches the landing; a citizen on the stair goes to the nearest landing and waits upon the Peacekeeper's word.
\item Every further order of the Mayor under this Article reaches a floor by the same road, and is posted at the landing with the hour, and no order binds a floor before it is posted there; a Peacekeeper who posts an order enters the hour, and the entries of every landing are gathered by the Deputy and delivered to the archive when the emergency is lifted.
\item Word that cannot wait for the stair passes by the stations as Article 21 allows, to the Deputy or Peacekeeper of the floor and no further; word from a floor to the Mayor passes by porter, or by the Peacekeeper to the Deputy's station and thence by the station of official word.
\end{clauses}

\section{The Sealing of a Floor}

A floor is sealed only against a danger that would otherwise pass beyond it, and a sealed floor is not a floor abandoned.

\begin{clauses}
\item The Mayor may order a floor, or a range of floors, sealed against the passage of citizens where the emergency is a disease under Article 19, a fouling of the air or the water, or a breach of the silo's structure; the order names the floors, the cause, the Peacekeepers who keep the seal at the landings above and below, and the day it is to be reviewed, which is not more than one cycle after; and no seal stands beyond one season save by the Assembly's renewal under Section 3.
\item A sealed floor draws the ration from its own station under Article 11, filled by Supply's runners to the landing and no further, and draws its water and air as before; where the emergency is a fouling of the water or the air, Mechanical carries the sealed floor's water to the landing and Supply its air in such vessels as it keeps for the purpose, and the Infirmary's aide is sent within the seal and remains there.
\item A citizen within a sealed floor may send word by the Peacekeeper at the landing, who carries it to the Deputy's station as Section 6 provides, and the household of a citizen sealed on a floor not the citizen's own is told where the citizen is; a citizen may not pass the seal in either direction save a Physician, a Peacekeeper, or a citizen of Mechanical the Head sends, and a citizen who does is answerable under Article 17 as for the violation of isolation under Article 19.
\item A citizen who dies within a sealed floor is examined by the Physician within the seal, and the Form of Death under Article 19 is sent out by the Peacekeeper, and the body is kept within the seal in such manner as the Physician orders until the seal is lifted, and given to the soil under Article 11 then.
\item The Judge reviews every seal at the day the order names, upon the reports of the Physician, the Head of Mechanical, and the Deputy, and enters whether the cause stands; a seal whose cause the Judge finds to have passed is lifted the same day, whatever the Mayor's order said.
\end{clauses}

\section{The Ration in Emergency}

The ration is lessened in emergency as it is lessened at any shortfall, alike for every citizen; this section sets down only the order in which the silo's other draws are lessened before it.

\begin{clauses}
\item Upon a shortfall of food, water, air, or power in an emergency, the Mayor orders the lessening in this order and no other, each draw lessened to what the emergency requires before the next is touched: the market under Article 12 is closed; the recreation of Article 19 and the lights of the common halls are stopped; the servers' share under Article 8 is lessened; the farms' share under Article 11 is lessened, and the beds let go dry in the order the Head of Supply enters; the reserve under Article 11 is opened by the Mayor's written order; and last, the ration itself is lessened under Article 11, alike for every citizen.
\item The works of Article 9, the Infirmary, and the stations of official word are not lessened in any emergency save where the Head of Mechanical enters that the generator cannot bear them, and then by the Head's order and not the Mayor's.
\item The Head of Supply reports to the Mayor and the Judge at each cycle of an emergency what the stores hold and how many cycles they will hold at the emergency's draw, and the report is posted at every landing with the Mayor's orders under Section 6.
\item A citizen who draws, keeps, or sells in an emergency beyond the ration the emergency allows is answerable under Article 17 as for the diversion of stores, and the grace of that Article's first offense is not given for it.
\end{clauses}

\section{The Labor of Emergency}

The silo's work does not stop for an emergency; it changes hands.

\begin{clauses}
\item The Mayor may assign any citizen of majority to any labor the emergency requires, beyond the citizen's trade and beyond the ordinary hours of Article 13, by an order posted at the landing naming the floor, the labor, and the Head under whom it is done; a citizen so assigned works under that Head as under the citizen's own, and the Head enters the citizen in the Department's roll for the emergency's span.
\item The rest day under Article 1 is not owed in an emergency, and no citizen is owed a substitute for it; but no citizen labors more than two shifts in any day save the tenders of the works, and every citizen assigned is fed at the labor's floor from the station there.
\item A citizen assigned to a critical role's labor under Article 13 in an emergency, or to the works, or to the workings, works only beside a citizen confirmed in that trade and does nothing alone; a citizen so assigned is not thereby of the trade, and is not vetted, and is bound by the entry Article 8, Section 11 provides where the labor was that Department's.
\item A citizen hurt or killed in the labor of emergency is entered in the roll of accidents of the Department under whose Head the labor was done, and the household petitions for restitution under Section 5 and Article 17; a citizen who cannot return to the citizen's own trade after is reassigned or retired under Article 13 as for a hurt in the trade.
\end{clauses}

\section{The Lifting of an Emergency}

An emergency ends as it began, by a declaration told to every floor, and what it cost is counted after.

\begin{clauses}
\item The Mayor lifts an emergency by a declaration sent and posted as Section 6 provides, naming the day and the hour; every seal, every assignment of labor, every lessening under Section 8, and every order given under this Article ceases at that hour, save such as the Judge continues under Section 3 for a stated cause and a stated span, entered.
\item Within one cycle of the lifting, every Department returns its citizens to their own trades and floors, Housing under Article 14 lodges anew any citizen whose dwelling the emergency took, and Supply resumes the ration in full and reports to the Mayor the day the reserve will be made good.
\item Within the same cycle, the Deputies deliver to the archive the entries of every landing under Section 6, the Physicians the Forms of every death within a seal, and every Head the roll of the emergency's labor and accidents; the Judge's review under Section 5 begins upon the delivery, and is read at the next sitting of the Assembly under Article 20.
\item The Mayor's report at that sitting states the emergency's cause, its span, the measures ordered, the stores spent, the citizens hurt and dead, and what the green list under Section 2 is to be changed to on the emergency's account; the change is entered before the next ordinary sitting, as Section 2 provides, and the Assembly is told that it has been.
\end{clauses}
